\begin{table*}[!ht]
    \setlength{\tabcolsep}{4pt}
    % \scriptsize
    \floatconts
      {tab:survival_3_6}
      {\caption{Prediction Performance Across Time Snapshots (3 and 6 months)}\vspace{-.75em}}
      {\adjustbox{max width=.9\textwidth}{\begin{tabular}{l  c c c @{\hspace{0.5cm}} c c c }
    \toprule
    \textbf{Snapshot Time} &\multicolumn{3}{c}{3 months} & \multicolumn{3}{c}{6 months} \\
    \midrule
     & \(C^{td}\) & IBS & \(\text{AUC}^{td}\) & \(C^{td}\) & IBS & \(\text{AUC}^{td}\) \\
    \midrule
    CoxPH & 0.639 ± 0.008 & 0.198 ± 0.004 & 0.681 ± 0.012 & 0.633 ± 0.011 & 0.189 ± 0.003 & 0.682 ± 0.011 \\
    Random Survival Forest & {\bftab 0.659} ± 0.008 & {\bftab 0.190} ± 0.003 & {\bftab 0.705} ± 0.010 & {\bftab 0.661} ± 0.012 & {\bftab 0.181} ± 0.004 & {\bftab 0.705} ± 0.011 \\
    DeepSurv & 0.646 ± 0.018 & 0.196 ± 0.005 & 0.689 ± 0.022 & 0.630 ± 0.010 & 0.190 ± 0.005 & 0.670 ± 0.010 \\
    DeepHit & 0.570 ± 0.036 & 0.240 ± 0.008 & 0.588 ± 0.044 & 0.575 ± 0.050 & 0.219 ± 0.018 & 0.613 ± 0.043 \\
    \bottomrule
    \end{tabular}}\vspace{-1em}}
\end{table*}


We assess the performance of survival models using three key evaluation metrics, computed on the test set across five independent experimental repeats. 
Each repeat involves a different randomized train-validation-test split (\(60/20/20\)). 
Reported values include the mean and standard deviation across repeats.

\paragraph{Survival Model Metrics.}  
The following metrics are used to evaluate time-to-event predictions:
\begin{enumerate}[label=(\roman*),itemsep=0pt]
    \item \textbf{Time-dependent Concordance Index (\(C^{td}\))} \citep{antolini2005time}:  
    Measures the model’s ability to correctly rank adverse event times over different time horizons.
    Higher values indicate better discriminative performance.
    
    \item \textbf{Integrated Brier Score (IBS)} \citep{graf1999assessment}:  
    Assesses the overall prediction accuracy survival probability by integrating the Brier score over time:
    \[
    \text{IBS} = \int_{0}^{M} \text{BS}(t)\, w(t)\, dt,
    \]
    where \(\text{BS}(t)\) is the Brier score at time \(t\) and \(w(t)\) is a weighting function. Lower values denote better calibration.
    
    \item \textbf{Time-dependent Area Under the Curve (\(\text{AUC}^{td}\))}: Quantifies the model’s discriminative ability by evaluating how well it distinguishes individuals experiencing events at different time points \citep{uno2007evaluating}. 
    We report the mean across all prediction time points. 
    A higher values suggest improved predictive performance.
\end{enumerate}

\paragraph{Causal Inference Evaluation.}  
At each snapshot \(\tau\), the ATE is estimated on the full cohort.
The survival models used for ATE estimation are trained on different randomized subsets across experimental repeats, employing cross-validation to avoid overfitting. 
The ATE is reported along with its standard deviation across experimental repeats. The standard deviation gives an empirical measure of uncertainty in the~ATE.

In our experimental setup, a negative ATE suggests that non-adherence is associated with an earlier occurrence of adverse events (in months), and a positive ATE implies a delay in adverse events occurring.